% Section structure adapted from 6.pdf; draft prompts are not thesis prose.
\chapter{GPU Index Design and Construction}
\label{ch:architecture}
A compressed GPU index needs both a searchable representation and a
construction process that fits device memory. This chapter develops
JHQ's training, encoding and list-ordered storage \cite{jhq}, then checks
the resulting representation. Query execution and measured construction
cost follow in \cref{ch:gpu-execution,sec:7-7}.

\section{Build and Search Workflow}
\label{sec:4-1}
Index construction separates reusable models from database-specific
arrays. Training fits the models; add encodes the database and places
each row in IVF-list order. Search then uses this completed, GPU-resident
representation.

\begin{figure}[!htbp]
\centering
\resizebox{\linewidth}{!}{% Build and search workflow: host lane above, GPU lane below, stages left to right.
\begin{tikzpicture}[
 font=\small\linespread{1}\selectfont,
 box/.style={draw=ThesisNavy!65,rounded corners=2pt,align=center,
   text width=3.25cm,minimum height=1.6cm,inner sep=4pt,fill=white},
 gpu/.style={box,fill=ThesisNavy!10},
 arr/.style={-{Latex[length=2mm]},draw=ThesisNavy!80,thick},
 lab/.style={font=\footnotesize,fill=white,inner sep=1.5pt},
 lane/.style={font=\small\bfseries,text=ThesisNavy,align=center}]
% lanes
\fill[ThesisNavy!3] (-3.2,1.15) rectangle (14.55,-1.6);
\fill[ThesisNavy!7] (-3.2,-1.6) rectangle (14.55,-4.85);
\draw[densely dashed,ThesisNavy!60] (-3.2,-1.6) -- (14.55,-1.6);
\node[lane] at (-2.45,-0.2) {Host\\(CPU)};
\node[lane] at (-2.45,-3.2) {GPU};
% stage headers
\node[font=\small\bfseries] at (0,1.75) {1\quad Train};
\node[font=\small\bfseries] at (6.0,1.75) {2\quad Build (\texttt{add})};
\node[font=\small\bfseries] at (12.1,1.75) {3\quad Search};
\draw[ThesisNavy!60] (-1.65,1.45) -- (1.65,1.45);
\draw[ThesisNavy!60] (2.45,1.45) -- (9.65,1.45);
\draw[ThesisNavy!60] (10.45,1.45) -- (13.75,1.45);
% host lane
\node[box] (hmodel) at (0,-0.2) {Samples; rotation $\Pi$\\Primary codebook\\Model cache};
\node[box] (hdata) at (4.05,-0.2) {Database vectors,\\read in bounded\\batches};
\node[box] (hquery) at (12.1,-0.2) {Query batch;\\returned IDs\\and scores};
% GPU lane
\node[gpu] (gtrain) at (0,-3.2) {Fit or load\\IVF centroids and\\residual codebooks};
\node[gpu] (gbuild) at (4.05,-3.2) {Assign and encode\\Order and pack\\(1 or 2 build passes)};
\node[gpu] (gindex) at (8.05,-3.2) {Resident index:\\offsets, IDs,\\codes, corrections};
\node[gpu] (gsearch) at (12.1,-3.2) {Search pipeline\\(Chapter~5)};
% GPU flow
\draw[arr] (gtrain) -- (gbuild);
\draw[arr] (gbuild) -- (gindex);
\draw[arr] (gindex) -- (gsearch);
% host <-> GPU transfers
\draw[arr] ([xshift=-8mm]hmodel.south) -- node[lab,pos=.45,left=1pt] {samples, models} ([xshift=-8mm]gtrain.north);
\draw[{Latex[length=2mm]}-{Latex[length=2mm]},draw=ThesisNavy!80,thick,dashed] ([xshift=8mm]gtrain.north) -- node[lab,pos=.55,right=1pt] {save/load} ([xshift=8mm]hmodel.south);
\draw[arr] (hdata) -- node[lab,pos=.45,right=1pt] {batches} (gbuild);
\draw[arr] ([xshift=-8mm]hquery.south) -- node[lab,pos=.45,left=1pt] {queries} ([xshift=-8mm]gsearch.north);
\draw[arr] ([xshift=8mm]gsearch.north) -- node[lab,pos=.55,right=1pt] {results} ([xshift=8mm]hquery.south);
\end{tikzpicture}
}
\caption[Build and search workflow]{Build and search workflow.}
\label{fig:4-architecture}
\end{figure}

The host prepares samples, the rotation matrix and the analytical primary
codebook; the GPU applies the rotation, fits or loads the IVF/residual
models, and encodes database batches into the resident index. Reusable
model state comprises the rotation, codebooks and centroids. The dashed
path shows saving and loading this state: a cache hit skips model fitting,
while add still encodes and orders database rows.

One successful \texttt{add} fixes the physical positions. Further adds
are rejected, and updates require rebuilding. Returned scores and the
host-call timing boundary are defined in
\cref{eq:5-refined-constant,sec:3-3}.


\section{Training, Rotation and Codebook Generation}
\label{sec:4-2}
\subsection{Transform and sampling}

QR factorisation on the host splits a seeded Gaussian matrix into
orthogonal and upper-triangular factors; its orthogonal factor gives
square $\Pi$. Training samples and database vectors use the same transform:
\begin{equation}
 Y_t=X_t\Pi^{\mathsf T},\qquad y=\Pi x,\qquad
 \|\Pi(q-x)\|_2^2=\|q-x\|_2^2
 \quad\text{when }\Pi^{\mathsf T}\Pi=I.
 \label{eq:4-transform}
\end{equation}
Dimension remains $d$. The validation in \cref{sec:4-6} distinguishes
stored-matrix geometry (K0) from rounding in general matrix
multiplication (GEMM).

The evaluated version (v69) stores FP32 matrices but enables cuBLAS
TF32 tensor-core math for $d\geq1{,}536$. This affects
training/add/query rotation and centroid GEMMs on both OpenAI datasets;
the implementation's environment switch cannot disable this mode. Thus FP32 storage does
not imply full FP32 multiplication. The FP32 variant (v72)
corrects this switch. Its sensitivity experiment reuses the same
TF32-trained cache and changes add/query arithmetic, rather than retraining
the model in FP32 (\cref{sec:7-precision}).

Primary/IVF training selects regularly spaced database rows, requesting
$n_t=39n_{\rm list}$:
\[
 n_t=\begin{cases}159{,}744,&\text{Vogue},\\319{,}488,&\text{arXiv/OpenAI},\\1{,}277{,}952,&\text{BGE/Stella}.\end{cases}
\]
This is Faiss's sample-count warning threshold, rather than a guarantee
of quantisation quality \cite{faiss_clustering}.

\subsection{Analytical primary codebook}

The primary codebook is generated from the sample scale. Its Cartesian
levels preserve the coordinate structure used by the encoder and lookup
tables:
\begin{equation}
 \sigma^2=\frac{1}{n_t d}\sum_{i=0}^{n_t-1}\|x_i\|_2^2,
 \qquad D_s=d/M,\qquad K=2^B.
 \label{eq:4-scale}
\end{equation}
\[
 B=D_s=8\Longrightarrow
 \mathcal C_m=\{-a,+a\}^{8},\quad a=0.674\sigma.
\]
JHQ encodes these eight signs in one byte in $O(d)$ work \cite{jhq}.
JHQ-GPU exploits this structure to split its lookup tables
(\cref{sec:5-2}). Unconstrained Lloyd updates can break this
Cartesian structure. For fixed sign vector $s\in\{-1,+1\}^d$, the
squared distance expands as
\begin{equation}
 \|q'-a s\|^2=\|q'\|^2+da^2-2a\langle q',s\rangle.
 \label{eq:4-sign-ranking}
\end{equation}
\[
 a>0\Longrightarrow
 \arg\min_s\|q'-as\|^2=\arg\max_s\langle q',s\rangle.
\]
For $a>0$, primary ranking therefore depends on the query/sign inner
product without requiring Gaussian data. Residuals still depend on $a$;
zero scale and FP32 ties need separate handling.

\subsection{Coarse assignment as filter and refine}

IVF initialises centroids from spaced rotated rows and runs eight
iterations using $s_\ell=\|\mu_\ell\|^2-2\langle y,\mu_\ell\rangle$.
Supported dimensions use int8 scores to filter candidates, then resolve
ambiguity in FP32. The following lemma states the condition under which
interval filtering preserves the assignment.

\begin{resultblock}
\begin{lemma}[Assignment with certified score intervals]
\label{lem:4-assignment-bound}
Let $\widetilde y,\widetilde\mu_\ell$ be dequantised approximations,
$e=y-\widetilde y$ and $f_\ell=\mu_\ell-\widetilde\mu_\ell$.
Keeping centroid norms exact, the approximate score
$\widetilde s_\ell=\|\mu_\ell\|^2-2\langle\widetilde y,\widetilde\mu_\ell\rangle$
has real-arithmetic error bounded by
\begin{equation}
 E_\ell=2\bigl(\|y\|\|f_\ell\|+\|e\|\|\widetilde\mu_\ell\|\bigr).
 \label{eq:4-assignment-bound}
\end{equation}
Given valid intervals $[L_\ell,U_\ell]$ containing every target score,
discarding only $L_\ell>\min_j U_j$ and exactly resolving the remaining
scores preserves the minimum and its specified tie rule.
\end{lemma}
\begin{proof}
Using $e=y-\widetilde y$ and $f_\ell=\mu_\ell-\widetilde\mu_\ell$,
\[
 \begin{aligned}
 \langle y,\mu_\ell\rangle-
 \langle\widetilde y,\widetilde\mu_\ell\rangle
 &=\langle y,f_\ell\rangle
   +\langle e,\widetilde\mu_\ell\rangle,\\
 |s_\ell-\widetilde s_\ell|
 &\leq2\bigl(\|y\|\|f_\ell\|
       +\|e\|\|\widetilde\mu_\ell\|\bigr)=E_\ell .
 \end{aligned}
\]
For $j^*\in\arg\min_j U_j$,
\[
 L_\ell>U_{j^*}
 \Longrightarrow
 s_\ell\geq L_\ell>U_{j^*}\geq s_{j^*}.
\]
The strict inequality excludes $\ell$; equality stays eligible.
\end{proof}
\end{resultblock}

The guarantee requires valid score intervals, including accumulation
error. Separate FP16 assignment checks found differences on BGE and
Stella; normal cold-build logs do not identify the assignment precision.
These checks therefore do not certify identical build assignments.

\subsection{Residual fitting and the iteration budget}

Residuals are defined relative to the primary reconstruction,
$r=\Pi x-\widehat p(x)$. Residual fitting uses 100,000 rows. Within each
subspace, GPU sorting and Lloyd updates fit $K_r=2^{B_r}$ sorted scalar
levels shared across coordinates, with a cap of 2,000 updates.

Earlier fixed-primary/IVF tests on Vogue and OpenAI-3072 support bounded
residual sampling: at 25 updates, using 100,000 rather than all rows
lost at most 0.0018 recall. Increasing updates from 25 to 2,000 gained
0.0063--0.0076 recall. Sample-size effects were non-monotone, and some
levels continued changing at the iteration cap. The cap bounds training
work; it does not establish convergence.

\subsection{Training work and list imbalance}

The training-work model separates transform generation, sample rotation,
primary-codebook generation and residual extraction. If $n_r$ denotes
the residual-training row count, the transform terms are
\[
 W_{\rm QR}=O(d^3),\qquad W_{\rm rotate}=O(n_t d^2).
\]
Direct coarse assignment costs $O(n_t n_{\rm list}d)$ per iteration;
the remaining terms are
\[
 W_{\rm primary}=O(MKD_s),\qquad
 W_{\rm residual\ extraction}=O(n_rd),
\]
followed by sorting and scalar updates. These are operation counts.
Synchronised diagnostics place about 83\% of BGE/Stella training in IVF
k-means (about 2.97\,s). They use one add buffer and forced
synchronisation, separately from the normal overlapped measurements in
\cref{sec:7-7}. Unequal list lengths affect subsequent scan work and
motivate \cref{sec:5-5}; the effects of eight IVF iterations, spaced seeds
and finite sampling on that imbalance were not measured separately.


\section{IVF Assignment, IDs and Physical Memory Layout}
\label{sec:4-3}
\subsection{From database IDs to list-ordered positions}

List ordering makes the vectors assigned to one IVF centroid contiguous.
An input row has identifier $i\in\{0,\ldots,N-1\}$; coarse assignment
uses computed centroid scores with the mathematical target
\begin{equation}
 a(i)=\operatorname*{arg\,min}_{0\leq\ell<n_{\rm list}}
       \|\Pi x_i-\mu_\ell\|_2^2.
 \label{eq:4-assignment}
\end{equation}
The implementation uses the numerical assignment path described above.
Primary and residual encoding do not subtract $\mu_{a(i)}$: list
membership restricts the search universe independently of those
reconstructions.

The builder sorts assignments together with input IDs. Let $\rho(p)$
be the input ID at sorted physical position $p$, and $o_\ell$ the start
of list $\ell$. The completed arrays satisfy
\begin{equation}
 0=o_0\leq o_1\leq\cdots\leq o_{n_{\rm list}}=N,
 \qquad p\in[o_\ell,o_{\ell+1})\Longrightarrow a(\rho(p))=\ell.
 \label{eq:4-offsets}
\end{equation}
Lower-bound searches find each list's first position; equal offsets
therefore mark an empty list:
\[
 o_\ell=o_{\ell+1}\Longleftrightarrow\text{list }\ell\text{ is empty}.
\]
Positions stay fixed in a completed index but can change after
rebuilding. All arrays use the same $\rho(p)$, and IDs map
$p\mapsto\rho(p)$ back to input rows. Two-pass add temporarily builds
$\rho^{-1}$ to write new encodings directly at their destination
(scatter). Primary ties use physical positions; final ties use database
IDs.

\par\medskip\noindent\begin{minipage}{\linewidth}
\captionsetup{hypcap=false}
\centering
% Independent arrays share a physical position; the band replaces connector lines.
\begin{tikzpicture}[x=1cm,y=1cm,font=\small\linespread{1}\selectfont,
 cell/.style={draw=ThesisNavy!45,minimum width=.86cm,minimum height=.48cm,inner sep=1pt},
 panel/.style={draw=ThesisNavy!35,rounded corners=2pt,fill=ThesisNavy!2}]
\draw[panel] (0,-1.85) rectangle (15.1,-.85);
\draw[panel] (0,-4.15) rectangle (15.1,-2.15);
\draw[panel] (0,-5.55) rectangle (15.1,-4.55);
\draw[panel] (0,-6.95) rectangle (15.1,-5.95);
\fill[ThesisNavy!12] (6.22,-6.90) rectangle (7.08,.3);
\node[font=\small\bfseries] at (6.65,1.12) {$p=3$};
\node[anchor=west,align=left] at (.15,.03) {$o=[0,3,3,8]$\\list 1 is empty};
\node[anchor=west,font=\small\bfseries] at (11.2,.03) {Byte offset};
\foreach \p in {0,...,7} {\node[cell] at (3.95+.90*\p,0) {$\p$};}
\draw[decorate,decoration={brace,amplitude=3pt}] (3.52,.32)--(6.18,.32) node[midway,above=4pt] {list 0};
\draw[decorate,decoration={brace,amplitude=3pt}] (6.22,.32)--(10.68,.32) node[midway,above=4pt] {list 2};
\node[anchor=west,align=left] at (.2,-1.35) {\textbf{Input IDs}\\$\mathrm{IDs}[p]=\rho(p)$};
\foreach \p/\id in {0/5,1/1,2/7,3/0,4/6,5/3,6/4,7/2} {\node[cell] at (3.95+.90*\p,-1.35) {$\id$};}
\node[anchor=west] at (11.2,-1.35) {$4p$};
\node[anchor=west,align=left] at (.2,-2.95) {\textbf{Primary words}\\4 codes per word};
\foreach \j/\yy in {0/-2.70,1/-3.30} {
 \foreach \p in {0,...,7} {\node[cell] at (3.95+.90*\p,\yy) {$\mathsf w_{\j,\p}$};}}
\node[anchor=west] at (11.2,-2.98) {$4(jN+p)$};
\node at (7.1,-3.83) {Adjacent candidates (threads) read consecutive words};
\node[anchor=west,align=left] at (.2,-5.05) {\textbf{Residual rows}\\$b_r$ bytes per row};
\foreach \p in {0,...,7} {\node[cell] at (3.95+.90*\p,-5.05) {$r_{\p}$};}
\node[anchor=west] at (11.2,-5.05) {$pb_r+t$};
\node[anchor=west,align=left] at (.2,-6.45) {\textbf{Corrections}\\FP32};
\foreach \p in {0,...,7} {\node[cell] at (3.95+.90*\p,-6.45) {$\gamma[\p]$};}
\node[anchor=west] at (11.2,-6.45) {$4p$};
\end{tikzpicture}

\captionof{figure}[Physical GPU index layout]{Physical GPU index layout.}
\label{fig:4-layout}
\end{minipage}\par\medskip

The example uses $N=8$ and offsets $o=[0,3,3,8]$, so list~1 is empty. The shaded column marks $p=3$: four separate arrays use the same physical position, and byte offsets are relative to each array's base. At fixed word group $j$, adjacent candidates occupy consecutive 32-bit words, allowing threads scanning adjacent candidates to coalesce their reads. In the eight-bit residual configuration, each row holds the $b_r$ code bytes of one vector.

\subsection{Exact addresses and representation widths}

The primary representation groups candidates by subspace. For primary
byte $Z[m,p]$, its address is
\[
 \operatorname{offset}_{\rm byte}(Z[m,p])=mN+p.
\]
Final packing combines four subspace bytes of one vector into each word:
\begin{equation}
 \mathsf w[j,p]=\sum_{t=0}^{3}2^{8t}Z[4j+t,p],\qquad
 \operatorname{addr}_{\rm byte}(\mathsf w[j,p])=4(jN+p).
 \label{eq:4-primary-address}
\end{equation}
Here $j=0,\ldots,M/4-1$. $\text{word-layout build}\Longrightarrow M\bmod4=0$;
other $M$ are rejected, while a separate byte-layout compilation retains
the byte representation. Each array occupies $NM$ bytes; repacking
temporarily holds both copies, costing $2NM$. The one-pass builder
reorders existing codes (gather) into a vector-major $[N,M]$ array,
transposes it and repacks, freeing obsolete copies. Two-pass construction
scatters directly into subspace-major storage before the same repack.
The search-time lane mapping is discussed in \Cref{sec:5-3}.

Residual codes remain row-major, with $b_r=\lceil dB_r/8\rceil$ bytes per
vector. In the eight-bit configuration, coordinate $t$ of row $p$ is
at byte offset $pb_r+t$. In the four-bit path it is in byte
$pb_r+\lfloor t/2\rfloor$, with even coordinates in the low four bits
(low nibble) and odd coordinates in the high four bits (high nibble).
A four-byte correction per physical row stores
\begin{equation}
 \gamma[ p ]=2\langle\widehat p(x_{\rho(p)}),
                              \widehat r(x_{\rho(p)})\rangle.
 \label{eq:4-correction}
\end{equation}
The correction must move with its row's codes during gathering or
scattering. List-order scans can then fetch residuals only for retained
candidates. Refinement uses compressed reconstructions; raw FP32
reranking would require an additional $4Nd$ device bytes. Chapter 7
evaluates the resulting recall.


\section{Resident Storage, Workspace and Index Lifecycle}
\label{sec:4-4}
\subsection{Stored bytes and allocated bytes}

Persistent row data are only part of device allocation: models, query
workspace and temporary copies are additional. For byte-sized primary
codes, row and list-offset storage are
\begin{equation}
 S_{\rm rows}=N\left(M+\left\lceil\frac{dB_r}{8}\right\rceil+8\right),
 \qquad S_{\rm offsets}=4(n_{\rm list}+1).
 \label{eq:4-row-storage}
\end{equation}
IDs and corrections contribute $8N$ bytes. With $M=d/8$, $B_r=8$,
the code payload alone is
\[
 b_{\rm code}=d/8+d=9d/8=1.125d.
\]
This exceeds 8-bit scalar quantisation (SQ8) by 12.5\%; graph storage
and vector-code bytes must therefore be considered together.

\begin{table}[!htbp]
\centering
\caption[Per-vector storage model]{Per-vector bytes. JHQ includes codes, ID and correction; SQ8 is
vector codes alone; CAGRA int8 models vector codes plus 32 four-byte
neighbours. Offsets, model parameters, workspace and other runtime
allocations are excluded. These components are not complete process totals.}
\label{tab:4-bytes}
\footnotesize\setlength{\tabcolsep}{4pt}
\begin{tabular}{@{}rrrrrrrr@{}}
\toprule
$d$ & Primary & Residual & ID + corr. & JHQ & SQ8 & \shortstack{CAGRA\\int8} & FP32 \\
\midrule
768 & 96 & 768 & 8 & 872 & 768 & 896 & 3,072 \\
1,024 & 128 & 1,024 & 8 & 1,160 & 1,024 & 1,152 & 4,096 \\
1,536 & 192 & 1,536 & 8 & 1,736 & 1,536 & 1,664 & 6,144 \\
3,072 & 384 & 3,072 & 8 & 3,464 & 3,072 & 3,200 & 12,288 \\
\bottomrule
\end{tabular}
\end{table}

At graph degree $d_g=32$, CAGRA int8 uses $d+4d_g$ row bytes. Stella's
JHQ row payload (19.20\,GiB) is close to int8 CAGRA's reported index
allocation of 19,656\,MiB, with different accounting boundaries.
Changing $\alpha$ changes residual reads, not the stored row bytes.

Basic model storage accounts for the rotation matrix, primary and
residual codebooks, and centroids with their norms:
\[
 S_{\rm model}
 =4d^2+4MKD_s+4MK_r+4n_{\rm list}(d+1)
 \quad\text{bytes}.
\]
Derived arrays are additional. Query workspace holds two query matrices,
centroid scores, primary LUTs, probes and retained candidates; refinement
scores residuals directly. The reported 128-probe configuration reserves
about 93.8\,MiB for scan groups even when unused
(\cref{eq:5-group-workspace}). Allocated capacity therefore also affects
resident memory.

The input remains caller-owned. For $n_b$ pipeline buffers with
$a_{\rm in}$ staged rows each, pinned host staging and device ordering
arrays use
\[
 S_{\rm pinned}=4n_ba_{\rm in}d,\qquad
 S_{\rm order}=S_{\rm inverse}=4N.
\]
The inverse map is temporary. Assignment scratch is freed before
gathering (\cref{eq:4-staging-rows,eq:4-assignment-scratch}), so host
staging, resident data and construction peaks have distinct lifetimes.

\subsection{Cache reuse and invalidation}

Cache keys omit code versions and full-data hashes, so code/data changes
require invalidation. Parameter sweeps reuse models; cold builds refit
them. Similar recall alone does not establish model identity. A cache
hit still incurs QR, model upload and add; warm-up and scheduler setup
are also needed before steady-state search. Construction and startup
timing boundaries are compared in \cref{sec:7-7}.

Changing $\alpha$ needs no retraining. Query buffers are reused by CUDA
Graph replay, and workspace grows before capture. With the result count
fixed, changes to batch size, survivors, probes, lanes, groups or bounds
trigger recapture. Calibration may bypass capture; index-owned buffers
remain allocated until release.


\section{Construction Cost and Tested Scale}
\label{sec:4-5}
\subsection{Bounded input batches and two-pass construction}

Keeping whole-database raw and rotated arrays on the GPU would require
$4Nd+4Nd=8Nd$ bytes. Bounded staging limits these arrays while the
compressed index remains resident. For $n_b$ pipeline buffers,
configured add rows $a_{\rm add}$ and per-array budget $S_{\rm stage}$,
the staged row count and raw/rotated device storage are
\begin{equation}
 a_{\rm in}=\max\!\left\{1,\min\!\left(a_{\rm add},
 \left\lfloor\frac{S_{\rm stage}}{4d}\right\rfloor\right)\right\},
 \qquad S_{\rm float}=8n_ba_{\rm in}d.
 \label{eq:4-staging-rows}
\end{equation}
\[
 4d\leq S_{\rm stage},\quad n_b=3,\ S_{\rm stage}=128\,\mathrm{MiB}
 \Longrightarrow S_{\rm float}\leq2n_bS_{\rm stage}=768\,\mathrm{MiB}.
\]
The bound covers these floating-point arrays; assignment and encoder
scratch are additional. One-pass construction encodes rows in input
order and gathers them into list order. Two-pass construction first
assigns and orders rows, then rereads and encodes directly into their
final positions, trading an extra traversal for fewer compressed copies.

\par\medskip\noindent\begin{minipage}{\linewidth}
\captionsetup{hypcap=false}
\centering
\begin{tikzpicture}[font=\small\linespread{1}\selectfont,
 stage/.style={draw=ThesisNavy!55,fill=ThesisNavy!4,rounded corners=2pt,
 text width=3.0cm,align=center,minimum height=12mm,inner sep=3pt},
 arrow/.style={-{Latex[length=2mm]},draw=ThesisNavy!70}]
\node[stage] (a) at (0,0) {Read and rotate};
\node[stage] (b) at (3.85,0) {Assign and encode\\Input order};
\node[stage,fill=ThesisNavy!12] (c) at (7.70,0) {Gather\\List-order copy};
\node[stage] (d) at (11.55,0) {Repack\\Primary codes};
\node[stage] (e) at (0,-2.3) {Assign raw rows\\Sort positions};
\node[stage] (f) at (3.85,-2.3) {Reread and rotate\\Encode};
\node[stage] (g) at (7.70,-2.3) {Write directly\\Into list order};
\node[stage,fill=ThesisNavy!12] (h) at (11.55,-2.3) {Repack\\Primary codes};
\foreach \x/\y in {a/b,b/c,c/d,e/f,f/g,g/h}\draw[arrow] (\x)--(\y);
\node[anchor=south west,font=\small\bfseries] at (-1.5,.85) {One pass};
\node[anchor=south west,font=\small\bfseries] at (-1.5,-1.45) {Two passes};
\node[align=center] at (7.70,-1.0) {Stella row arrays: 40.59 GiB};
\node[align=center] at (11.55,-3.3) {Stella row arrays: 21.32 GiB};
\end{tikzpicture}

\captionof{figure}[One- and two-pass construction]{One- and two-pass construction.}
\label{fig:2-add-paths}
\end{minipage}\par\medskip

For Stella, analytical row-array occupancy is 40.59\,GiB for one-pass gather and 21.32\,GiB for two-pass repack (\cref{tab:4-phase-storage}). One pass holds both input-order and list-order encodings; two passes write directly into list order. In the first assignment pass, centroids are transformed back with $\Pi^\top$, avoiding rotation of database rows.

\subsection{Memory use across construction phases}

The implementation estimates one-pass memory from the row arrays and
an allowance for other allocations:
\begin{equation}
 H_{\rm single}=N(2b_r+3M+20)+2^{30}\quad\text{bytes}.
 \label{eq:4-peak-heuristic}
\end{equation}
The ten live arrays contribute
\[
 3NM+2Nb_r+20N=N(3M+2b_r+20).
\]
The 1\,GiB allowance is heuristic. During final two-pass repacking,
IDs, residuals, corrections and two primary copies instead require
\begin{equation}
 S_{\rm repack,2}=N(b_r+2M+8).
 \label{eq:4-two-pass-repack}
\end{equation}
Whole-build peak also includes models, sorting, context, earlier
pipeline/inverse scratch and later query workspace. The phase values in
\Cref{tab:4-phase-storage} isolate row arrays, rather than whole-process
peaks.

\begin{table}[!htbp]
\centering
\caption[Analytical row-array occupancy by construction phase]{Analytical row-array occupancy (GiB), $d=1{,}024$, $M=128$,
$B_r=8$. BGE uses one pass; Stella uses two. The one-pass Stella column
is counterfactual capacity analysis. Pipeline, model, offsets and runtime
costs are excluded from these phase values.}
\label{tab:4-phase-storage}
\small
\begin{tabular}{@{}lrrr@{}}
\toprule
Dataset & Final row payload & One-pass gather & Two-pass repack \\
\midrule
BGE & 10.90 & 23.05 & 12.11 \\
Stella & 19.20 & 40.59 & 21.32 \\
\bottomrule
\end{tabular}
\end{table}

Assignment-product buffers have a separate capacity $a_{\rm assign}$,
configurable independently of staged input. The evaluated default is 32,768
for at least 8,192 lists, otherwise 8,192. Their storage is
\begin{equation}
 S_{\rm dots}=4n_b n_{\rm list}a_{\rm assign}\quad\text{bytes}.
 \label{eq:4-assignment-scratch}
\end{equation}
Allocation capacity is independent of the active/staged rows. At
$d=1536/3072$, both defaults allocate 3\,GiB. For BGE/Stella,
$3\times4=12$\,GiB assignment scratch is released before gathering.
For BGE, the 12\,GiB scratch and the 23.05\,GiB row-array gather therefore do not coexist. The peak is
the largest sum of allocations live at the same time:
\[
 S_{\rm peak}=\max_t\sum_{A\ {\rm live\ at}\ t}|A|.
\]



Stella's cold-build log selects two passes: the one-pass estimate in
\cref{eq:4-peak-heuristic} is 41.6\,GiB, versus the log's reported
30.6\,GiB tail-stage availability. Separately sampled resident process
totals, including context and query workspace, are 11.93\,GiB for BGE
and 20.44\,GiB for Stella (\cref{sec:4-4}). Supplementary 20\,ms device
samples observe a 26.08\,GiB BGE peak near the start of add and a
22.58\,GiB Stella peak during repack. The smaller two-pass repack is
consistent with Stella's lower observed peak. These sampled device
totals include context and can miss shorter-lived peaks.


\section{Index Representation and Validation}
\label{sec:4-6}
Validation must check the representation independently of successful
searches. A signed-32-bit $n_td$ overflow at $2^{31}$ previously set the
codebook scale to zero, and a cache without version checks reloaded it.
Chunked FP64 norm accumulation fixed the fault; routing success alone
had not exposed it.

\subsection{Representation consistency checks}

Pilot exports cover Vogue, arXiv and OpenAI-3072. The exporter shares
the evaluated implementation's encoding, training and construction code;
this source comparison covers the shared path, while trained-model
identity and large-dataset two-pass construction need separate checks.
The first checks therefore examine the exported representation rather
than re-encoding original input vectors.

Codebooks passed separability, scalar-level and sorting checks. An FP64
loop checked $\gamma=2\langle\widehat p,\widehat r\rangle$ on every exported
row, with maximum absolute error $1.28\times10^{-7}$--$2.77\times10^{-7}$.
Sampled GPU-returned top-10 scores have maximum relative errors from
$2.38\times10^{-7}$ to $4.85\times10^{-5}$, using
$|s_{\rm GPU}-s_{64}|/s_{64}$ with the query-norm constant included.
These top-10 returned-score checks skip IDs absent from the routed map, so
they do not fully check output membership.

The pilot also checked proposed residual-distance bounds against FP64
reconstructed distances across seven workloads, 896 query instances
and 72,800,590 routed candidate instances, with zero reported violations.
The same candidates support the FP64 composite-score identity check;
these representation checks are separate from the selection tests below.

\subsection{Input encoding and data placement checks}
\label{sec:4-representation-audit}

The supplementary input audit extends to Vogue, arXiv, OpenAI-3072 and
Stella. Offsets are monotone with endpoints 0 and $N$, and IDs form a
permutation of all database rows in each dataset. FP64 products of the
stored rotation give maximum entrywise $\|\Pi\Pi^{\mathsf T}-I\|$ errors
of $3.96$--$4.34\times10^{-7}$. Sampled norm checks use FP64 rotation
rather than GPU GEMM, measuring the stored matrix's geometry.
\Cref{tab:4-representation-audit} summarises recoding from original inputs.

\par\medskip\noindent\begin{minipage}{\linewidth}
\centering\captionsetup{hypcap=false}
\captionof{table}[Original-input FP64 recoding audit]{Original-input FP64 recoding audit. Flips count coordinate
decisions; the final column is the largest residual-boundary distance
among mismatches. Structural checks cover all rows; recoding is sampled.}
\label{tab:4-representation-audit}
\scriptsize\setlength{\tabcolsep}{3pt}
\begin{tabular}{@{}lrrrrr@{}}
\toprule
Dataset & Rows & Orth. error & Primary flips & Residual flips & Max margin \\
\midrule
Vogue & 200,000 & 4.11e-07 & 14 & 5,409 & 1.12e-07 \\
arXiv & 100,000 & 4.11e-07 & 1 & 2,515 & 7.31e-08 \\
OpenAI-3072 & 20,000 & 3.96e-07 & 5,745 & 1,646,806 & 2.67e-05 \\
Stella & 100,000 & 4.34e-07 & 16 & 3,064 & 1.03e-07 \\
\bottomrule
\end{tabular}

\end{minipage}\par\medskip

At 768/1,024 dimensions, recoding mismatches lie within
$1.12\times10^{-7}$ of decision boundaries. OpenAI-3072 has 2.68\%
residual mismatches, up to $2.67\times10^{-5}$ from boundaries,
consistent with TF32. Residual recoding uses stored primary digits so
that primary flips are not counted again as residual errors.

Comparing one-pass and two-pass builds by original ID gives identical
codes on the three smaller databases, corrections within
$8.94\times10^{-8}$, and 1/1/40 Vogue/arXiv/OpenAI assignment differences.
Stella checks cover only two-pass construction. Exported centroids are
missing, so list membership remains unverified.

\subsection{Validation coverage and limitations}
\suppressfloats[t]

\par\medskip\noindent\begin{minipage}{\linewidth}
\captionsetup{hypcap=false}
\centering
\captionof{table}[Representation contracts and audit coverage]{Representation contracts after the pilot and supplementary input audit.
Checks separate geometry, arithmetic and row placement.}
\label{tab:4-contracts}
\footnotesize\setlength{\tabcolsep}{4pt}
\begin{tabularx}{\linewidth}{@{}>{\raggedright\arraybackslash}p{1.8cm}*{2}{>{\raggedright\arraybackslash}X}@{}}
\toprule
Boundary & Required property & Evidence / remaining check \\
\midrule
K0:\newline geometry & Common transform with bounded floating-point isometry error & Stored-matrix orthogonality checked on four datasets; TF32 GPU arithmetic distinguished \\
K1:\newline codebooks / correction & Cartesian levels, sorted residual codebooks, correction matches decoded codes & Pilot checks completed on three databases; all-row correction errors in \cref{sec:4-6} \\
K1:\newline input encoding & Stored codes encode the intended original input row & FP64 original-input recoding on four samples; boundary and TF32 discrepancies quantified \\
K1:\newline placement & Offsets cover rows; IDs form a permutation; all fields follow it & All-row offsets/ID checks; three ID-joined path comparisons; centroid membership unverified \\
Execution & Valid indexing, synchronisation and score-order contracts & Archived query and sanitizer checks; finite coverage, detailed in \cref{sec:7-4} \\
\bottomrule
\end{tabularx}
\end{minipage}\par\medskip

The evaluated packed implementation requires
\[
 M\mid d,\quad D_s\mid B,\quad M\bmod4=0.
\]
Inputs must be finite and counts valid; unused slots need sentinel
values marking missing entries. IDs are signed 32-bit and address
products use wider arithmetic. The strict selection check validates candidate selection
(\cref{sec:7-strict-selection}); centroid membership, empty-list
behaviour and adversarial tails remain uncertified.

